\documentclass[10pt,a4paper]{amsart}
\usepackage{amsmath,amssymb,amsthm,graphicx}
\usepackage{verbatim}
\newenvironment{reference}{\comment}{\endcomment}
%\newenvironment{reference}{}{}
\newenvironment{slogan}{\comment}{\endcomment}
\newenvironment{history}{\comment}{\endcomment}

% For commutative diagrams we use Xy-pic
\usepackage[all]{xy}

% We use 2cell for 2-commutative diagrams.
\xyoption{2cell}
% Theorem environments.
%
\theoremstyle{plain}
\newtheorem{theorem}[subsection]{Theorem}
\newtheorem{proposition}[subsection]{Proposition}
\newtheorem{lemma}[subsection]{Lemma}

\theoremstyle{definition}
\newtheorem{definition}[subsection]{Definition}
\newtheorem{example}[subsection]{Example}
\newtheorem{exercise}[subsection]{Exercise}
\newtheorem{situation}[subsection]{Situation}

\theoremstyle{remark}
\newtheorem{remark}[subsection]{Remark}
\newtheorem{remarks}[subsection]{Remarks}

\numberwithin{equation}{subsection}

% Macros
%
\def\lim{\mathop{\mathrm{lim}}\nolimits}
\def\colim{\mathop{\mathrm{colim}}\nolimits}
\def\Spec{\mathop{\mathrm{Spec}}}
\def\Hom{\mathop{\mathrm{Hom}}\nolimits}
\def\Ext{\mathop{\mathrm{Ext}}\nolimits}
\def\SheafHom{\mathop{\mathcal{H}\!\mathit{om}}\nolimits}
\def\SheafExt{\mathop{\mathcal{E}\!\mathit{xt}}\nolimits}
\def\Sch{\mathit{Sch}}
\def\Mor{\mathop{\mathrm{Mor}}\nolimits}
\def\Ob{\mathop{\mathrm{Ob}}\nolimits}
\def\Sh{\mathop{\mathit{Sh}}\nolimits}
\def\NL{\mathop{N\!L}\nolimits}
\def\CH{\mathop{\mathrm{CH}}\nolimits}
\def\proetale{{pro\text{-}\acute{e}tale}}
\def\etale{{\acute{e}tale}}
\def\QCoh{\mathit{QCoh}}
\def\Ker{\mathop{\mathrm{Ker}}}
\def\Im{\mathop{\mathrm{Im}}}
\def\Coker{\mathop{\mathrm{Coker}}}
\def\Coim{\mathop{\mathrm{Coim}}}

% Boxtimes
%
\DeclareMathSymbol{\boxtimes}{\mathbin}{AMSa}{"02}

%
% Macros for moduli stacks/spaces
%
\def\QCohstack{\mathcal{QC}\!\mathit{oh}}
\def\Cohstack{\mathcal{C}\!\mathit{oh}}
\def\Spacesstack{\mathcal{S}\!\mathit{paces}}
\def\Quotfunctor{\mathrm{Quot}}
\def\Hilbfunctor{\mathrm{Hilb}}
\def\Curvesstack{\mathcal{C}\!\mathit{urves}}
\def\Polarizedstack{\mathcal{P}\!\mathit{olarized}}
\def\Complexesstack{\mathcal{C}\!\mathit{omplexes}}
% \Pic is the operator that assigns to X its picard group, usage \Pic(X)
% \Picardstack_{X/B} denotes the Picard stack of X over B
% \Picardfunctor_{X/B} denotes the Picard functor of X over B
\def\Pic{\mathop{\mathrm{Pic}}\nolimits}
\def\Picardstack{\mathcal{P}\!\mathit{ic}}
\def\Picardfunctor{\mathrm{Pic}}
\def\Deformationcategory{\mathcal{D}\!\mathit{ef}}
\usepackage{cleveref}
\crefname{originalreference}{original reference}{original references}
\Crefname{originalreference}{Original reference}{Original references}
\makeatletter
\newlabel{situation-shriek}{{0F42}{source}{}{sourcepointer-1}{}}
\newlabel{situation-shriek@cref}{{[originalreference][0][]0F42}{[1][0][]source}{}{}{}}
\makeatother
\crefname{originalreference}{original reference}{original references}
\Crefname{originalreference}{Original reference}{Original references}
\makeatletter
\newlabel{remark-local-calculation-shriek}{{0BV2}{source}{}{sourcepointer-2}{}}
\newlabel{remark-local-calculation-shriek@cref}{{[originalreference][0][]0BV2}{[1][0][]source}{}{}{}}
\makeatother
\crefname{originallemma}{original lemma}{original lemmas}
\Crefname{originallemma}{Original lemma}{Original lemmas}
\makeatletter
\newlabel{dualizing-lemma-relative-dualizing-algebraic}{{0BZW}{source}{}{sourcepointer-3}{}}
\newlabel{dualizing-lemma-relative-dualizing-algebraic@cref}{{[originallemma][0][]0BZW}{[1][0][]source}{}{}{}}
\makeatother
\crefname{originallemma}{original lemma}{original lemmas}
\Crefname{originallemma}{Original lemma}{Original lemmas}
\makeatletter
\newlabel{perfect-lemma-affine-locally-rel-perfect}{{0DI2}{source}{}{sourcepointer-4}{}}
\newlabel{perfect-lemma-affine-locally-rel-perfect@cref}{{[originallemma][0][]0DI2}{[1][0][]source}{}{}{}}
\makeatother
\title{The shriek of the structure sheaf along a flat morphism is relatively perfect and self-dual}
\author{The Stacks Project Authors}
\date{stacks-project:duality.tex:local-snapshot-20261012; GNU FDL 1.2 or later}
\begin{document}
\maketitle
\noindent\textit{Editorial note.} All delivered mathematical text is the author's own, copied from the pinned Stacks Project chapter file duality.tex; no statement or proof step was rewritten, reordered or removed.\par\medskip
\noindent\textit{Editorial note.} Preamble provenance: the pinned chapter file opens with an input of the shared preamble and carries no preamble text of its own, so the project's single Stacks-wide preamble staged at sources/stacks-duality/preamble.tex is excerpted verbatim under the recipe fixed by LANE\_BRIEF 22.1, as three spans, lines 8-13 (verbatim package and the reference/slogan/history environments), lines 14-18 (Xy-pic setup) and lines 154-217 (the native theorem declarations and mathematical macros). Nothing was added, removed or redefined; the wrapper class and the amsmath/amssymb/amsthm/graphicx base packages are supplied by the assembler.\par\medskip
\noindent\textit{Editorial note.} The counted statement and proof contain no cross reference, no citation and no included file, so the delivered body is self contained.\par\medskip
\noindent\textit{Editorial note.} Canonical location of the counted proposition: Stacks tag 0E9T (sources/stacks-duality/pinned-original-source.tex line 4890).\par\medskip
\noindent\textit{Allowed standard prerequisites (editorial).}\begin{itemize}
\item The local calculation of the shriek on affines. Same chapter, Remark labelled remark-local-calculation-shriek, cited at proof lines 4901-4904. Turns the statement into an algebraic one.
\item The algebraic statement for a flat ring map. Dualizing Complexes, Lemma labelled dualizing-lemma-relative-dualizing-algebraic, cited at proof lines 4904-4906. Supplies both conclusions in the affine case.
\item Comparison of relatively perfect complexes with the algebraic notion. Derived Categories of Schemes, Lemma labelled perfect-lemma-affine-locally-rel-perfect, cited at proof lines 4904-4906. Matches relative perfection with the algebraic hypothesis.
\end{itemize}\medskip
\ifdefined\hypertarget\hypertarget{sourcepointer-1}{}\fi
\noindent\textit{Source reference.} Original reference 0F42: sources/stacks-duality/pinned-original-source.tex, source lines 3953--3953. Same-chapter pointer into the author environment situation at source line 3953; invoked at proof line None of the counted proof. It is retained only as a navigation pointer; every load-bearing step is inside the delivered span.\par\medskip
\ifdefined\hypertarget\hypertarget{sourcepointer-2}{}\fi
\noindent\textit{Source reference.} Original reference 0BV2: sources/stacks-duality/pinned-original-source.tex, source lines 4636--4636. Same-chapter pointer into the author environment remark at source line 4636; invoked at proof line 4901 of the counted proof. It is retained only as a navigation pointer; every load-bearing step is inside the delivered span.\par\medskip
\ifdefined\hypertarget\hypertarget{sourcepointer-3}{}\fi
\noindent\textit{Source reference.} Original lemma 0BZW: sources/stacks-chow/tags.txt, source lines 15066--15066. Cross-chapter standard prerequisite located in the project tag list staged in the local snapshot; invoked at proof line 4903 of the counted proof. It is retained only as a navigation pointer; every load-bearing step is inside the delivered span.\par\medskip
\ifdefined\hypertarget\hypertarget{sourcepointer-4}{}\fi
\noindent\textit{Source reference.} Original lemma 0DI2: sources/stacks-chow/tags.txt, source lines 16927--16927. Cross-chapter standard prerequisite located in the project tag list staged in the local snapshot; invoked at proof line 4905 of the counted proof. It is retained only as a navigation pointer; every load-bearing step is inside the delivered span.\par\medskip
\begin{lemma}
\label{lemma-flat-shriek-relatively-perfect}
In Situation \ref{situation-shriek} let $f : X \to Y$ be a morphism of
$\textit{FTS}_S$. If $f$ is flat, then
$f^!\mathcal{O}_Y$ is a $Y$-perfect object of $D(\mathcal{O}_X)$ and
$\mathcal{O}_X \to
R\SheafHom_{\mathcal{O}_X}(f^!\mathcal{O}_Y, f^!\mathcal{O}_Y)$
is an isomorphism.
\end{lemma}

\begin{proof}
Both assertions are local on $X$. Thus we may assume $X$ and $Y$ are
affine. Then Remark \ref{remark-local-calculation-shriek}
turns the lemma into an algebra lemma, namely
Dualizing Complexes, Lemma \ref{dualizing-lemma-relative-dualizing-algebraic}.
(Use Derived Categories of Schemes, Lemma
\ref{perfect-lemma-affine-locally-rel-perfect} to match the languages.)
\end{proof}

\end{document}
